% Paper09b v6 (2026-10-17), layout-ready successor to v5; v5 is preserved.
% Equation/table typography only; source data, verifier, thresholds, and conclusions are unchanged.
\documentclass[10pt,twocolumn]{article}
\usepackage[a4paper,margin=15mm]{geometry}
\usepackage{amsmath,amssymb}
\usepackage{amsthm}
\usepackage{graphicx}
\graphicspath{{paper09b_arv_certification_v6_src/}{../}{./}}
\theoremstyle{plain}
\newtheorem{theorem}{Theorem}
\newtheorem{lemma}[theorem]{Lemma}
\newtheorem{proposition}[theorem]{Proposition}
\newtheorem{corollary}[theorem]{Corollary}
\theoremstyle{remark}
\newtheorem{remark}[theorem]{Remark}
\usepackage{booktabs,array,calc}
\usepackage[font=small,labelfont=bf]{caption}
\usepackage[colorlinks=true]{hyperref}
\providecommand{\tightlist}{\setlength{\itemsep}{0pt}\setlength{\parskip}{0pt}}
\emergencystretch=3em
\allowdisplaybreaks

\begin{document}
\title{Part II --- Conditional Regime Viability on the Northern Cod Stock:\\
Exact Accounting Thresholds on Public Assessment and Landings Data}
\author{Amin Abaee\\[0.35em]
{\small Independent Researcher}\\[0.55em]
{\small\href{https://orcid.org/0000-0002-0019-1842}{ORCID: 0000-0002-0019-1842}}\\[0.3em]
{\small\href{mailto:amin\_abaee@ut.ac.ir}{amin\_abaee@ut.ac.ir}}}
\date{October 17, 2026}
\maketitle

\begin{abstract}
This paper develops a conditional accounting analysis of public Northern cod (NAFO 2J3KL) records, keeping the 1983--2015 NCAM input vintage and the 1954--2024 xteNCAM reassessment distinct. The assessed quantity is spawning-stock biomass (SSB), \(B_t\), in kilotonnes; reported catches and landings are separate series, not SSB-equivalent removals. DFO 2025 Table~1 reports total landings in tonnes, not all fishing removals. Since \(1{,}000\) tonnes equal \(1\) kilotonne, write \(C_y^{\rm rep}\) for reported landings/catch in kt and \(K_y=\kappa C_y^{\rm rep}\) for unknown SSB-equivalent removals, with \(\kappa\ge0\). Source timing does not identify whether a year's landings align with transition \(t\to t+1\) as \(C_t^{\rm rep}\) or \(C_{t+1}^{\rm rep}\); both assignments are evaluated.

For a declining SSB transition, the implied accounting multiplier is subunitary under either of two placements if and only if \(\kappa<\kappa^*=(B_t-B_{t+1})/C_y^{\rm rep}\), for the selected reported-series row. For 1991--1992 the thresholds are \(2.0496\) and \(8.6083\); for 1992--1993 they are \(6.8586\) and \(24.6577\). Across both steps, assignments, and placements, the conditional contraction result therefore holds only for \(0\le\kappa<2.0496\). The post-collapse thresholds and a separate 2020--2021 database-vintage comparison are likewise conditional. Because \(\kappa\) and within-year alignment are unresolved, the result is not an unconditional biological or causal conclusion. Survey, assessment-band, and threshold-crossing analyses describe the record but do not identify the conversion or a mechanism. Exact-rational checks reproduce the stated threshold arithmetic on the printed inputs; they do not validate the source estimates.
\end{abstract}

\textbf{Keywords:} viability; obstruction; regime shift; Northern cod;
exact computation; monitoring design; limit reference points

\section{Introduction}\label{intro}

\noindent\textbf{Companion relation.} This Part II accompanies
\emph{Certification, not simulation: the viability of harvest rules on two real resource systems,
and the horizon of what can be certified} (paper09 main). It shares the Northern cod
study-system record named there but not the main paper's fitted-model method:
the interval-specific inputs, exact-rational arithmetic, and brackets required for
this document's certificates are stated here so that it compiles and verifies independently.



The viability and obstruction machinery used here --- survivable
sets, robust multipliers, breach-time certificates, and probe design ---
is developed in Abaee (2026, An obstruction calculus) and validated on
small exact instances. This paper applies the accounting component to
the Northern cod stock (NAFO divisions 2J3KL). Its empirical issue is
one of commensurability: an assessed SSB change and a reported landings
row are not interchangeable quantities, and the source record does not
fix a unique within-year alignment. The paper therefore derives
row-specific \(\kappa^*\) thresholds under two explicit accounting
placements and two adjacent-year assignments. It does not identify
actual SSB-equivalent removals, reconstruct a biological zero-harvest
trajectory, or adjudicate the historical cause of the collapse. The
forecast layer of the same stock record --- competing forecast ladders
and the 2022--2023 observation suspension as a review-timing event ---
is analyzed separately in the companion paper (Abaee, 2026, The
Northern cod forecast ladder, Section 3.10). The two papers share a
system record but make distinct inferential claims.

The claims below are of three kinds, and the label travels with the
claim. \emph{Arithmetic claims} are exact identities or inequalities on
the displayed inputs; the accompanying checker regenerates the
row-specific \(\kappa^*\) values and tests the two accounting forms,
while other descriptive comparisons are stated with their source rows.
\emph{Provenance claims} identify file, vintage, and column, including
whether a quantity is a model input or output. \emph{Scope claims} are
conditional readings --- stress bounds, sustained hypotheticals, or
diagnostic comparisons --- and are labelled where they appear. No
interpretive claim is presented as a source identity. The contributions
are: the assessment-record spine and its provenance ledger (Sections
\ref{spine}--\ref{windows}); the row-specific \(\kappa^*\) thresholds,
which state exactly when a specified accounting multiplier would be
subunitary (Section~\ref{brackets}); and the window profile and
assessment-band summaries (Section~\ref{profile}). The threshold
results are conditional on the unknown conversion from reported
landings/catch to SSB-equivalent removals and identify no biological
mechanism.

\section{Related work}
\label{priorart}

This paper derives row-specific thresholds for when a non-fishing
accounting factor would be subunitary on particular Northern cod
assessment-SSB steps, conditional on an unknown conversion from reported
landings/catch to SSB-equivalent removals and on either of two stated
accounting placements. It does not estimate that conversion or identify
a biological cause. This distinction matters when comparing the result
with the empirical literature on the collapse.

\paragraph{The received attribution.} Hutchings and Myers (1994) concluded that the collapse
of northern cod could be attributed solely to overexploitation. Myers and Cadigan (1995)
tested the competing hypothesis that an increase in natural mortality in the first half of
1991 caused the collapse, did not support that hypothesis, and concluded that overfishing was
sufficiently high to cause the collapse of the population even allowing for the possibility
of elevated recent natural mortality. Myers, Hutchings and Barrowman (1997) rejected the
hypothesis that mortality of young cod is unrelated to fully recruited fishing mortality,
finding high juvenile mortality associated with high adult mortality in each of six stocks.

\paragraph{Why the distinction matters.} Those papers address the
\emph{historical cause} of the collapse. This paper does not adjudicate
that question. For an observed decline in the published SSB series, a
specified \(\kappa\), a selected reported-landings row, and either
simplified accounting placement, the algebra gives a row-specific
threshold \(\kappa^*\). Only when \(\kappa<\kappa^*\) does that
accounting imply a subunitary non-fishing factor. Because the reported
landings are not themselves SSB-equivalent removals, \(\kappa\) is
unknown, and the within-year alignment is unresolved, the result is a
conditional statement about the \emph{arithmetic of the assessment
record}, not a biological zero-harvest counterfactual or a causal
finding. Neither the reported landings series nor the SSB decline alone
identifies productivity, natural mortality, or the contribution of
fishing to the collapse.

\paragraph{What the accounting cannot resolve.} Myers and Cadigan
(1995) specifically declined to support elevated natural mortality as
the explanation. A conditional \(\kappa<\kappa^*\) result here must not
be read as evidence of a productivity collapse or as a refutation of
that finding: the conversion is not estimated, the accounting factors
are not biological state-transition parameters, and the data do not
identify a mechanism. Any such causal or biological interpretation
would go beyond the certificate.

\paragraph{Which series the certificates are computed on.} The concern is concrete rather
than hypothetical. Myers, Hutchings and Barrowman (1997) documented that for these stocks the
VPA-based and survey-based abundance trends diverge --- from the early 1980s the VPA-based
trend shows a decline where none is apparent in the survey-based trend --- and attributed the
divergence plausibly to discarding of young fish. Any certificate computed on one series
inherits that series' artefacts, so the choice of series is load-bearing for this paper and
cannot be left implicit.

This paper therefore checks the raw survey index directly. The survey
series also declines over 1989--1994 (Proposition~\ref{prop:survey}),
by the exact factor \(\tfrac{21797}{2127417}\), and its recorded
proportional decline is 2.6 times the assessment-series decline over
the common window. This is a comparison of two recorded series, not an
independent estimate of SSB or removals. The \(\kappa^*\) thresholds
are computed only from assessment SSB and reported landings/catch; the
survey index has different units and does not identify the conversion
or independently certify a non-fishing factor.

\paragraph{Position within viability theory.} That the governing question is which states
admit controls keeping every constraint satisfied, rather than which policy optimises, is the
founding move of viability theory (Aubin, 1991); the obstruction calculus used here is the
observation-constrained development of that programme (Abaee, 2026), and its constructive
counterpart on this same record is certified in the separate two-system paper09 main (System I). What this paper
contributes to that lineage is the \emph{conditional direction} of the
accounting thresholds --- when a specified conversion would imply a
subunitary factor under the declared forms --- and the fact that the
thresholds are obtained in exact rational arithmetic without fitting a
biological model.

\section{Data and provenance}\label{data}

Seven locked public files are used, each cross-verified against its
primary source at archival time (Abaee, 2026, Calibration data record): (i) the NCAM (Northern Cod Assessment Model) 2016
assessment input series (DFO, 2016, Table A2): spawning stock biomass
in kilotonnes, 1983--2015, cross-checked 33 of 33 years against the
primary report --- the \emph{input vintage} of the biomass record;
(ii) the xteNCAM 2025 reassessment SSB series (Regular et al., 2025,
Table 17), 1954--2024, with 95\% bands and SSB-over-\(B_{\mathrm{lim}}\)
ratios --- the \emph{reassessment vintage}; (iii) the reported total
landings series (Regular et al., 2025, Table 1), 1954--2023, in tonnes;
the source labels these values ``total landings'', not all fishing
removals; (iv) the catch reconstruction (Schijns et al., 2021, Table 1),
a separate reported series; (v) the fall research-vessel survey
abundance index (Schijns et al., 2021, Table 3, citing DFO 2021b) --- a
research-vessel input to the assessments, not an output of them; (vi)
the RAM Legacy database (Ricard et al., 2012) per-year extract for this
stock (the free 2024 release), 1850--2021 --- the \emph{database
vintage}, extending the series to 2021; and (vii) the capelin acoustic
biomass series (observed years only, no interpolation).

The assessment SSB \(B_t\) is expressed in kilotonnes (kt). The DFO
Table~1 landings column is reported in tonnes (t); the unit conversion is
\(1\,\mathrm{kt}=1{,}000\,\mathrm{t}\), so the analysis writes
\(C_y^{\rm rep}=L_y^{\rm rep}/1{,}000\) in kt. This conversion changes
units only; it does not make landings equivalent to SSB loss or to all
fishing removals. Let \(\kappa\ge0\), unknown and not estimated here,
map the reported series to SSB-equivalent removals:
\(K_y=\kappa C_y^{\rm rep}\) kt. For each assessed transition
\(t\to t+1\), the source record does not establish a unique within-year
placement; both \(C_t^{\rm rep}\) and \(C_{t+1}^{\rm rep}\) are therefore
carried through the threshold calculation.

Two thresholds appear throughout, and their provenance is stated once
here. The \emph{auxiliary floor} \(B_{\mathrm{aux}} = 276\) kt is the
reassessment vintage's \(B_{\mathrm{lim}}\) at the file's rounding: the
published ratio column is rounded to two decimals, and the values it
implies for \(B_{\mathrm{lim}}\) range from \(250\) to \(289\) kt
across the file's years. Every floor comparison below is insensitive
across that range (the compared quantities clear it by large factors;
Remark~\ref{rem:xtevintage}), so a single representative value is
fixed. The \emph{reference point} \(B_{\mathrm{ref}} = 884.6\) kt is
the input vintage era's limit reference point, whose documented
construction is the 1983--1989 mean of the assessment input series
(DFO, 2016); the exact mean is
\(\tfrac{44229}{50} = 884.58\) kt (Proposition~\ref{prop:windows}),
and \(884.6\) is its published rounding.

Two discipline rules are respected throughout. First, the mortality
and fishing-mortality columns of the assessment input file are
diagnostics only (dynamics are never built on them). Second, every
observed multiplier \(\rho_t=B_{t+1}/B_t\) is a \emph{net assessment-SSB}
quantity, not a non-fishing growth rate. Section~\ref{brackets}
compares it with the two declared accounting forms after introducing
\(K_y=\kappa C_y^{\rm rep}\); this is a conditional accounting
calculation, not a biological model.

\section{The record spine}\label{spine}

Table~\ref{tab:spine} tabulates the assessment-input SSB and reported
total-landings rows for the collapse era and its surroundings, alongside
the survey index, the covariate (observed years only), and the mortality
diagnostic. The landings column remains a reported series; the
SSB-equivalent quantity used in the conditional accounting is
\(\kappa C_y^{\rm rep}\), not the printed landings themselves.

\begin{table}[t]
\centering
\scriptsize
\setlength{\tabcolsep}{2.8pt}
\caption{The record spine, 1983--1995. SSB in kt (DFO 2016 assessment
input, Table A2); reported total landings in t (DFO 2025, Table 1,
which labels them ``total landings''; the Schijns et al. catch
reconstruction has identical reported rows for 1992--1995). These
landings are not labelled as all fishing removals or SSB-equivalent
losses. The survey index is in the research-vessel file's index units;
the covariate is in kt for observed years only; \(M\) is the
assessment's mean natural-mortality diagnostic (model-estimated).
Computed entries are exact rational arithmetic on the printed rows.}
\label{tab:spine}
\begin{tabular}{@{}r r r r r r@{}}
\toprule
Year & SSB (kt) & Landings (t) & Survey & Covariate & \(M\) \\
\midrule
1983 & 841.08 & 232{,}345 & 2{,}088{,}958 & --- & 0.391 \\
1984 & 863.23 & 232{,}471 & 2{,}198{,}605 & --- & 0.377 \\
1985 & 902.94 & 231{,}293 & 1{,}288{,}360 & 3426 & 0.349 \\
1986 & 836.00 & 266{,}713 & 2{,}502{,}702 & 3697 & 0.277 \\
1987 & 940.75 & 239{,}924 & 1{,}020{,}462 & 2576 & 0.494 \\
1988 & 886.40 & 268{,}677 & 1{,}223{,}314 & 4285 & 0.335 \\
1989 & 921.66 & 253{,}990 & 2{,}127{,}417 & 3712 & 0.289 \\
1990 & 861.92 & 219{,}452 & 1{,}627{,}647 & 5783 & 0.403 \\
1991 & 734.51 & 172{,}012 & 1{,}117{,}670 & 138 & 1.002 \\
1992 & 381.95 & 40{,}956 & 239{,}740 & 138 & 2.214 \\
1993 & 101.05 & 11{,}392 & 90{,}709 & --- & 2.575 \\
1994 & 30.55 & 1{,}314 & 21{,}797 & --- & 2.331 \\
1995 & 9.68 & 413 & 43{,}240 & --- & 0.288 \\
\bottomrule
\end{tabular}
\end{table}

\section{The reference window and the collapse window}\label{windows}

\begin{proposition}[seven readings, six transitions; the windows'
worst multipliers]
\label{prop:windows}
From the Table A2 series: the reference window 1983--1989 contains
seven readings \((841.08, 863.23, 902.94, 836.00, 940.75, 886.40,
921.66)\) kt and six annual transitions; their minimum is the 1985
transition \(\rho_{\mathrm{obs}} = 836.00/902.94 =
\tfrac{41800}{45147} \approx 0.9259 < 1\), and their exact mean is
\(\tfrac{44229}{50} = 884.58\) kt. The collapse window 1991--1993
contains two transitions,
\(\tfrac{38195}{73451} \approx 0.520\) (1992) and
\(\tfrac{2021}{7639} \approx 0.265\) (1993), each with numerator
strictly below denominator. The connecting 1990 reading is
\(861.92\) kt, with transitions
\(\tfrac{86192}{92166} \approx 0.935\) (1990) and
\(\tfrac{73451}{86192} \approx 0.852\) (1991).
\end{proposition}

\begin{proof}
Each multiplier is the exact ratio of consecutive Table A2 SSB entries.
The minimum of the six reference-window ratios and the two
collapse-window ratios follows by direct comparison of the printed
fractions. Each displayed fraction is below \(1\) because its
numerator is smaller than its denominator (\(41800 < 45147\),
\(38195 < 73451\), \(2021 < 7639\)). The mean is the exact rational
sum of the seven readings divided by seven.
\end{proof}

\section{Conditional non-fishing accounting thresholds}\label{brackets}

The observed multiplier \(\rho_t=B_{t+1}/B_t\) is a net assessment-SSB ratio. It folds together recruitment, natural mortality, fishing, and assessment estimation; it is not itself a non-fishing growth factor. To examine what a reported catch row could imply under simplified accounting, let \(C_y^{\rm rep}\) denote reported landings/catch in kt after the stated tonne-to-kiloton conversion, and let \(K_y=\kappa C_y^{\rm rep}\) be its SSB-equivalent counterpart, with the unknown \(\kappa\ge0\). No value of \(\kappa\) is estimated. For transition \(t\to t+1\), the source record does not establish whether the aligned annual row is \(C_t^{\rm rep}\) or \(C_{t+1}^{\rm rep}\); both are carried below. This year-alignment uncertainty is distinct from the two accounting placements, before or after the interval-wide factor.

\medskip
\noindent\textbf{Shared symbols across the companion papers.} The companion papers share this vocabulary with paper-local meanings; the letters below carry meanings here that differ there:
\begin{center}\small
\begin{tabular}{@{}p{0.40\linewidth}p{0.52\linewidth}@{}}
\toprule
\textbf{Here} & \textbf{In the companions} \\
\midrule
\(\rho_t=B_{t+1}/B_t\): the net assessment-SSB multiplier &
the mesh relaxation parameter (certification); the contamination radius
\(\rho\) of \(\mathcal{D}_{\rho}\) (sufficiency; minimax duals) \\
\(r_t=\kappa C_y^{\rm rep}/B_t\): the SSB-equivalent-removal share &
the intrinsic growth rates \(r\), \(r_i\) (forecast ladder; the
minimax benchmark) \\
\(B_t\): assessment spawning-stock biomass &
the belief/information sets \(B_t\), \(B\) (obstruction calculus;
certification; sufficiency) \\
\bottomrule
\end{tabular}
\end{center}

\begin{lemma}[conditional non-fishing multiplier bracket]\label{lem:bracket}
Fix a transition with assessment readings \(B_t,B_{t+1}>0\), a reported annual value \(C_y^{\rm rep}>0\) in kt assigned to that transition, and \(K_y=\kappa C_y^{\rm rep}<B_t\). Write \(\rho_t=B_{t+1}/B_t\), \(r_t=K_y/B_t\). Suppose the accounting uses one of two placements, with \(g_t>0\) the implied non-fishing accounting factor:
\[
\begin{aligned}
B_{t+1} &= g_t B_t-K_y && \text{(after-growth placement)},\\
B_{t+1} &= g_t(B_t-K_y) && \text{(before-growth placement)}.
\end{aligned}
\]
Then respectively \(g_t=\rho_t+r_t\) or \(g_t=\rho_t/(1-r_t)\), and
\[
\rho_t\le g_t\le\max\left\{\rho_t+r_t,\frac{\rho_t}{1-r_t}\right\}.
\]
For a declining assessment step \(B_t>B_{t+1}\), either placement gives \(g_t<1\) if and only if
\[
0\le\kappa<\kappa^*_{t,y}:=\frac{B_t-B_{t+1}}{C_y^{\rm rep}}.
\]
At equality \(g_t=1\); for \(\kappa\ge\kappa^*_{t,y}\), the step does not certify a subunitary factor under that alignment. If \(B_{t+1}\ge B_t\), no nonnegative \(\kappa\) gives a strict contraction in either form.
\end{lemma}

\begin{proof}
The after-growth equation gives \(g_t=(B_{t+1}+K_y)/B_t=\rho_t+r_t\). The before-growth equation gives \(g_t=B_{t+1}/(B_t-K_y)=\rho_t/(1-r_t)\), well-defined on the stated domain. Both are at least \(\rho_t\). For either form, \(g_t<1\) is equivalent to \(B_{t+1}+\kappa C_y^{\rm rep}<B_t\), or \(\kappa<(B_t-B_{t+1})/C_y^{\rm rep}\); equality yields one. The unitless threshold follows because both numerator and denominator are expressed in kt.
\end{proof}

The threshold is row-specific. If the source landings are retained in tonnes as \(L_y^{\rm rep}\), then \(C_y^{\rm rep}=L_y^{\rm rep}/1{,}000\) kt and
\[
\kappa^*_{t,y}=\frac{B_t-B_{t+1}}{C_y^{\rm rep}}=\frac{1{,}000(B_t-B_{t+1})}{L_y^{\rm rep}}.
\]
For a fixed transition, a conditional certificate that is robust to both adjacent-year assignments requires \(\kappa<\min\{\kappa^*_{t,t},\kappa^*_{t,t+1}\}\). The same strict condition covers both accounting placements. This is an exact implication of the declared equations, not a claim that the source identifies the true timing or \(\kappa\).

\begin{proposition}[row-specific thresholds for the selected transitions]\label{prop:brackets}
Using the printed assessment-input SSB readings and reported total-landings rows in Table~\ref{tab:brackets}, the collapse transitions 1991--1992 have thresholds \(2.0496\) for \(C_{1991}^{\rm rep}\) and \(8.6083\) for \(C_{1992}^{\rm rep}\); 1992--1993 have \(6.8586\) for \(C_{1992}^{\rm rep}\) and \(24.6577\) for \(C_{1993}^{\rm rep}\). Thus both collapse steps have subunitary accounting factors under both placements and both adjacent-year assignments only if \(\kappa<2.0496\). For the two later steps 1993--1994 and 1994--1995, the corresponding threshold pairs are \(6.1886,53.6530\) and \(15.8828,50.5327\); the paired post-collapse certificate therefore requires \(\kappa<6.1886\). The selected reference/prelude steps have joint alignment thresholds \(0.2510\) (1985--1986), \(0.2352\) (1989--1990), and \(0.5806\) (1990--1991). Because \(\kappa\) is not identified, none of these is an unconditional non-fishing or biological verdict.
\end{proposition}

\begin{proof}
For each table row and each adjacent-year landings placement, compute \(\kappa^*_{t,y}\) as the exact ratio of the printed SSB loss in kt to the converted reported-landings value in kt. The joint threshold is the smaller of the two row-specific ratios. The displayed decimals are rounded; the inequality is strict at the exact rational threshold.
\end{proof}

\begin{table*}[t]
\centering
\scriptsize
\setlength{\tabcolsep}{4pt}
\caption{Conditional thresholds for the selected assessment-SSB transitions. Reported total landings from DFO 2025 Table~1 are converted from tonnes to kt. For an assigned row \(C_y^{\rm rep}\), a subunitary accounting factor under either before- or after-growth placement follows only for \(\kappa<\kappa^*_{t,y}\). To be robust to both year alignments, \(\kappa\) must be below both thresholds in the last two columns.}
\label{tab:brackets}
\begin{tabular}{@{}lrrrrr@{}}
\toprule
Step & SSB loss (kt) & \(C_t^{\rm rep}\) (kt) & \(C_{t+1}^{\rm rep}\) (kt) & \(\kappa^*_{t,t}\) & \(\kappa^*_{t,t+1}\) \\
\midrule
1985$\to$86 & 66.94 & 231.293 & 266.713 & 0.2894 & 0.2510 \\
1989$\to$90 & 59.74 & 253.990 & 219.452 & 0.2352 & 0.2722 \\
1990$\to$91 & 127.41 & 219.452 & 172.012 & 0.5806 & 0.7407 \\
1991$\to$92 & 352.56 & 172.012 & 40.956 & 2.0496 & 8.6083 \\
1992$\to$93 & 280.90 & 40.956 & 11.392 & 6.8586 & 24.6577 \\
1993$\to$94 & 70.50 & 11.392 & 1.314 & 6.1886 & 53.6530 \\
1994$\to$95 & 20.87 & 1.314 & 0.413 & 15.8828 & 50.5327 \\
\bottomrule
\end{tabular}
\end{table*}

\section{The window profile and the band certificates}\label{profile}

The reference window's profile answers the selection question at all
horizons, and the reassessment vintage's published bands answer the
robustness question the point estimates leave open.

\begin{proposition}[the reference window's profile: declines only at
horizons one to three]\label{prop:profile}
Within the reference window 1983--1989, the minimum end-to-start ratio
of a \(k\)-year subwindow is: \(k=1\):
\(\tfrac{41800}{45147} \approx 0.9259\) (the worst step, at 2 of the 6
subwindows below unity); \(k=2\): \(\tfrac{83600}{86323} \approx
0.9685\) (2 of 5 below); \(k=3\): \(\tfrac{44320}{45147} \approx
0.9817\) (2 of 4 below); \(k = 4, 5, 6\): every subwindow grows ---
the minimum products are \(\tfrac{15361}{15049} \approx 1.0207\),
\(\tfrac{22160}{21027} \approx 1.0539\), and
\(\tfrac{15361}{14018} \approx 1.0958\). Fifteen of the twenty-one
subwindows grow in net from start to end; only five grow at
every constituent annual step. The sustained-decline
reading of the reference window therefore holds only at horizons one
to three; at every horizon of four or more years inside the window,
the record grew.
\end{proposition}

\begin{proof}
For each fixed \(k\), the subwindow products are exact ratios
\(S_{t+k}/S_t\) of the printed SSB values; their minima follow by
cross-multiplication. The \(k \ge 4\) growth statements are the direct
comparisons \(S_{t+k} > S_t\), which involve no roots. The subwindow
count is \(6+5+4+3+2+1 = 21\).
\end{proof}

\begin{proposition}[band certificates: exactly the collapse-era steps
certify, and the 1993 breach holds at the band's upper edge]
\label{prop:bands}
The reassessment vintage publishes 95\% bands \((B_{\mathrm{lo}},
B_{\mathrm{hi}})\) per year. A band-certified decline is a step whose
band interval contains only values below unity:
\([\,\tfrac{B_{\mathrm{lo}}(t+1)}{B_{\mathrm{hi}}(t)},
\tfrac{B_{\mathrm{hi}}(t+1)}{B_{\mathrm{lo}}(t)}\,] \subset (0,1)\).
Over the vintage's 1954--2024 record, the band-certified declines are
exactly the four collapse-era steps: 1991--1992
(\([\tfrac{265}{754}, \tfrac{468}{478}]\), upper bound \(0.9791\)),
1992--1993 (\([\tfrac{54}{468}, \tfrac{115}{265}]\), upper bound
\(0.4340\)), 1993--1994, and 1994--1995 --- and no other step in
seventy years, and no reference-decade step. The 1993 floor breach is
band-certified (\(B_{\mathrm{hi}}(1993) = 115 < 276\)); the 1992
reading is indeterminate on the bands (\(265 \le 276 \le 468\)).
\end{proposition}

\begin{proof}
For each adjacent pair the decline criterion reduces to
\(B_{\mathrm{hi}}(t+1) < B_{\mathrm{lo}}(t)\). The companion checker
compares all seventy consecutive-year pairs in the locked reassessment
extract using exact rational arithmetic and returns exactly the four
steps listed. The displayed interval endpoints are the corresponding
integer ratios, and the floor comparisons are direct rational comparisons.
\end{proof}

\begin{remark}[the vintage dependence of the typology]
\label{rem:vintagetypology}
The reassessment vintage reads the reference decade itself as
all-growth (its 1983--1989 multipliers are \(1.14, 1.22, 1.03, 1.28,
1.10, 1.01\), minimum \(1.01\)), with the decline beginning at
1989--1990 (\(738/864 \approx 0.854\)): the worst-step typology of
Proposition~\ref{prop:brackets} is vintage-internal, and the
paper's decade statements are made in the input vintage, whose
worst step (\(0.9259\)) has no counterpart in the reassessment
vintage. A third vintage (the database extract's assessment, 2021)
reads 1990--1993 as \(765, 747, 354, 83\) kt; all three vintages
agree on the breach sign pattern against the floor (above, above,
below).
\end{remark}

\begin{remark}[the record before the window: 1954--1982]
\label{rem:backward}
The reassessment vintage extends the record to 1954: a peak of
\(1508\) kt at 1962, with the 1983 window-opening reading at
\(427\) kt --- \(28.3\%\) of the peak, vintage-internally. The floor
\(276\) kt has a prior episode: the crossings of the point series are
exactly \(\{1977, 1980, 1993, 2016\}\) (below 1977--1979, recovered
1980--1992, below again 1993--2015, re-crossed 2016), and the 1978
reading is band-certified below the floor. The reference window thus
opens three years after a prior recovery, at \(28.3\%\) of the
record's peak scale: the window's start is a data artifact of the
input vintage's table, and the design critique of
Proposition~\ref{prop:discrimination} applies to the window's
placement as well as its mean.
\end{remark}

\section{Collapse-window landings and conditional post-collapse thresholds}\label{moratorium}

\begin{proposition}[conditional thresholds for the collapse-window steps]
\label{prop:moratorium}
Across 1991--1992 the assessment-input SSB series falls from \(734.51\) to \(381.95\) kt, a recorded decline of \(352.56\) kt. The reported landings are \(C_{1991}^{\rm rep}=172.012\) kt and \(C_{1992}^{\rm rep}=40.956\) kt, giving \(\kappa^*\) values \(2.0496\) and \(8.6083\), respectively. Under either accounting placement, a subunitary factor for this step is implied under both year assignments only for \(0\le\kappa<2.0496\). Across 1992--1993 the recorded SSB decline is \(381.95-101.05=280.90\) kt; the reported landings are \(40.956\) kt for 1992 and \(11.392\) kt for 1993, giving thresholds \(6.8586\) and \(24.6577\). Both year assignments imply a subunitary factor in either placement only for \(0\le\kappa<6.8586\). The two-step conditional statement therefore requires \(\kappa<2.0496\). Since neither \(\kappa\) nor within-year alignment is identified, the observed decline is not an unconditional non-fishing, biological, or causal certificate.
\end{proposition}

\begin{proof}
Each threshold is the assessment-SSB loss divided by the relevant annual reported-landings value after converting tonnes to kt. The exact ratios are those in Table~\ref{tab:brackets}; the joint conditions use the smaller threshold for each pair of adjacent rows. The SSB differences themselves are exact differences of the printed assessment-input values.
\end{proof}

\begin{proposition}[reported catch reconstruction and landings rows]
\label{prop:catchrecord}
The Schijns et al. catch-reconstruction rows for 1992--1995 are \(40{,}956, 11{,}392, 1{,}314, 413\) t, numerically identical to DFO 2025 Table~1's reported total-landings rows for those years; the 2015 catch-reconstruction row (\(4{,}436\) t) also matches the reported Table~1 value and the assessment report. These are cross-file equalities of reported series, not evidence that either source lists every fishing removal or gives SSB-equivalent losses. As a scale comparison only, each of the 1994--1995 reported-landings values is about 176--651 times smaller than the seven 1983--1989 values; this compares reported values, not SSB-equivalent removals.
\end{proposition}

\begin{proof}
The 1992--1995 and 2015 equalities are direct comparisons of the printed rows. Across all seven 1983--1989 Table~1 rows, the reported-landings range is \(231{,}293\)--\(268{,}677\) t; the two 1994--1995 rows are \(1{,}314\) and \(413\) t. Pairwise ratios therefore range from \(231{,}293/1{,}314\approx176.0\) to \(268{,}677/413\approx650.5\). This calculation does not convert either series into SSB-equivalent removals.
\end{proof}

\begin{proposition}[conditional thresholds for the post-collapse steps]
\label{prop:experiment}
For 1993--1994, the assessment-input SSB falls from \(101.05\) to \(30.55\) kt; the adjacent reported landings are \(C_{1993}^{\rm rep}=11.392\) kt and \(C_{1994}^{\rm rep}=1.314\) kt, with thresholds \(6.1886\) and \(53.6530\). For 1994--1995, the SSB falls from \(30.55\) to \(9.68\) kt; \(C_{1994}^{\rm rep}=1.314\) kt and \(C_{1995}^{\rm rep}=0.413\) kt, with thresholds \(15.8828\) and \(50.5327\). Thus both post-collapse steps have subunitary accounting factors under both placements and both year alignments only if \(0\le\kappa<6.1886\). The observed assessment SSB point series falls \(101.05\to30.55\to9.68\) kt, a ratio of \(\tfrac{10105}{968}\approx10.44\); that recorded change is distinct from the conditional accounting statement. The \(9.68\)-kt value is the minimum of the 1983--2015 Table A2 input series.
\end{proposition}

\begin{proof}
The four row-specific thresholds are the exact rational ratios in Table~\ref{tab:brackets}; the joint condition is their minimum across the two annual alignments and two steps. The two-year SSB ratio and point-series minimum follow from direct comparisons of the 33 printed Table A2 readings.
\end{proof}

\section{Mechanized verification}
\label{mech}

The two accounting identities in Lemma~\ref{lem:bracket} are formalized in the project's Lean~4 layer, in namespace \texttt{Formalizations.ARV} (\texttt{ARV\_RegimeViability}). The formal statements prove the ordered-field algebra for a nonnegative SSB-equivalent quantity \(K\) under the declared before- and after-growth forms, including the equivalence between a subunitary factor and the strict threshold inequality. They do not estimate \(\kappa\), validate the landings-to-SSB conversion, identify the within-year alignment, or verify the assessment inputs.

\noindent\textbf{Scope.} The Lean layer proves only the ordered-field algebra for a supplied nonnegative SSB-equivalent quantity \(K_y\) under the two declared accounting equations. It does not formalize the empirical map \(K_y=\kappa C_y^{\rm rep}\), the value of \(\kappa\), source extraction, unit conversion, row alignment, or assessment uncertainty. The accompanying Python checker addresses the selected source-row and threshold arithmetic, but does not re-extract the cited reports. Thus the formal result is: \emph{if} the supplied SSB readings, \(K_y<B_t\), and accounting placement satisfy the stated assumptions, the corresponding algebraic implication follows; no observation, biological mechanism, or causal explanation is machine-certified.

\section{Breach certificates and the cross-vintage check}\label{breach}

\begin{proposition}[breach certificates on the two thresholds]
\label{prop:breach}
With \(B_{\mathrm{aux}} = 276\) kt and \(B_{\mathrm{ref}} = 884.6\)
kt: (i) \emph{conditional stress bound}: from the end-1991 reading
\(734.51\) kt, one application of the observed 1992--1993 net SSB
multiplier already lands below \(B_{\mathrm{aux}}\), since
\(\tfrac{2021}{7639} < \tfrac{27600}{73451}\). This is a hypothetical
stress bound on the recorded net multiplier, not on \(g_t\) or on a
non-fishing factor; the realized assessment-series breach took two steps.
(ii) \emph{first below-floor assessment reading}: the 1992 reading is
above the floor (\(381.95 \ge 276\)) and the 1993 reading is below it
(\(101.05 < 276\)): the first under-floor point estimate in the input
vintage is 1993, and that vintage then stays below the floor through
2014 (\(250.12\) kt), re-crossing at 2015
(Proposition~\ref{prop:recovery}). (iii) \emph{reference point}:
\(734.51 < 884.6\): the 1991 reading is already below
\(B_{\mathrm{ref}}\), as is the 1990 reading before it
(\(861.92 < 884.6\)).
\end{proposition}

\begin{proof}
(i) is the exact cross-multiplication
\(2021 \cdot 73451 = 148{,}444{,}471 < 210{,}836{,}400 = 7639 \cdot
27600\). (ii) is the pair of direct rational comparisons
on the Table A2 rows, stated so that the intermediate year is visible:
the input-vintage point estimate first falls below the floor in 1993 only because its 1992 point estimate is above the floor. (iii) is \(734.51 < 884.6\) and \(861.92 < 884.6\).
No step of the proof uses floating point.
\end{proof}

\begin{remark}[the breach frontier]\label{rem:frontier}
On the realized path from the 1989 reading --- \(921.66 \to 861.92 \to
734.51 \to 381.95 \to 101.05 \to 30.55 \to 9.68\) kt --- the
first-below reading is a function of the threshold \(L\), certified on
intervals: \(L \in (861.92, 921.66]\) first falls below at the 1990
reading; \((734.51, 861.92]\) at 1991; \((381.95, 734.51]\) at 1992;
\((101.05, 381.95]\) at 1993; \((30.55, 101.05]\) at 1994;
\((9.68, 30.55]\) at 1995; and \(L \le 9.68\) is never breached within
1983--2015 (the series minimum is \(9.68\)). The realized threshold
\(B_{\mathrm{aux}} = 276\) sits in the 1993 interval
(Proposition~\ref{prop:breach}(ii)); the reference point
\(B_{\mathrm{ref}} = 884.6\) sits in the 1990 interval, after breaching
repeatedly within its own construction window
(Proposition~\ref{prop:discrimination}).
\end{remark}

\begin{remark}[the cross-vintage check]\label{rem:xtevintage}
The reassessment vintage reads the same three years as
\(600, 352, 79\) kt, with floor ratios \(2.18, 1.28, 0.29\). The two
vintages agree on the sign pattern against the floor --- above, above,
below --- and disagree on levels by certified factors
\(\tfrac{73451}{60000} \approx 1.22\),
\(\tfrac{38195}{35200} \approx 1.09\), and
\(\tfrac{10105}{7900} = \tfrac{2021}{1580} \approx 1.28\): the
reassessment's 1993 level is \(21.8\%\) below the input vintage's, and
its \(0.29\) ratio is its own level (\(79\) kt) over its own
\(B_{\mathrm{lim}}\)-scale denominator, not the input vintage's
\(101.05\). Because the ratio column is rounded to two decimals, the
implied \(B_{\mathrm{lim}}\) ranges over \(250\)--\(289\) kt across
the file; the sign pattern would flip only if the true threshold were
at or below \(101.05\) kt or above \(352\) kt, far outside that range,
so every floor comparison here is insensitive to the rounding. The
cross-file statement certified is therefore agreement on the breach
sign pattern, not equality of levels --- and the level disagreement
itself (\(\le\) a factor of \(1.28\) between vintages six years
apart) is part of the certified record.
\end{remark}

\section{The instruments}\label{instruments}

\begin{proposition}[the collapse is present in the raw survey index]
\label{prop:survey}
The fall research-vessel survey index reads \(2{,}127{,}417\) (1989),
\(1{,}117{,}670\) (1991), \(239{,}740\) (1992), \(90{,}709\) (1993),
and \(21{,}797\) (1994): a certified factor of
\(\tfrac{21797}{2127417} \approx 0.0102\) over 1989--1994, and
\(\tfrac{239740}{1117670} \approx 0.214\) across 1991--1992. The index
is a research-vessel input to the assessments, not an output of them:
the collapse is visible in the raw instrument record, before any model
processing. Over the common window 1989--1993 the index falls by
\(\tfrac{90709}{2127417} \approx 0.0426\) while the assessment series
falls by \(\tfrac{10105}{92166} \approx 0.1096\): the instrument falls
\(2.57\) times further than the model output it feeds.
\end{proposition}

\begin{proof}
All five values are rows of the locked survey file; the three factors
are their exact ratios, and \(2.57\) is the exact ratio
\(\tfrac{10105 \cdot 2127417}{92166 \cdot 90709}\) of the two
common-window factors. The provenance statement is a property of the
file (survey rows are assessment inputs, recorded at archival time).
The index-to-biomass mapping assumes constant catchability and
coverage; the divergence factor is a difference between the two
series' rates, not an attribution.
\end{proof}

\begin{proposition}[the covariate crossing is contemporaneous with the
mortality crossing]\label{prop:capelin}
In the observed capelin series, the acoustic estimate falls from
\(5783\) kt (1990) to \(138\) kt (1991) --- a certified factor of
\(\tfrac{138}{5783} \approx 0.024\) --- and reads \(138\) kt again in
1992. The Table A2 mean natural-mortality diagnostic crosses
\(M = 1\) across the same annual step (\(0.403\) in 1990 to \(1.002\)
in 1991; diagnostic column, model-estimated). At the data's annual
resolution the two crossings coincide; the covariate collapse precedes
the certified floor breach (the 1993 reading) by two annual steps and
the assessment series' steepest realized step (1991--1992) by one.
\end{proposition}

\begin{proof}
Both covariate readings are observed rows of the locked capelin file
(no interpolation, per the archival discipline); the factor is their
exact ratio. The mortality crossing is the diagnostic column's
consecutive values straddling \(1\) across the same pair of years as
the covariate collapse --- the rows display the coincidence directly,
and no sub-annual timing is claimed. Because the diagnostic is
model-estimated given assumed catch, its crossing corroborates the
regime reading but is not an independent certification of it. The
covariate is a single acoustic instrument; no comparability guarantee
across the early-1990s redistribution years is claimed, and the
certified statements are about the locked rows at annual resolution.
\end{proof}

\section{The thresholds as design objects}\label{design}

\begin{proposition}[the auxiliary floor separates the eras; the
reference point does not]\label{prop:discrimination}
Against \(B_{\mathrm{aux}} = 276\) kt: all seven reference-window
readings lie above the floor (minimum \(836.00\)), and the 1993--1995
readings (\(101.05, 30.55, 9.68\)) lie below it, with 1990--1992
above: the auxiliary floor separates the reference window from the
collapse era without exception. Against
\(B_{\mathrm{ref}} = 884.6\) kt: three of the seven reference-window
readings (\(841.08, 863.23, 836.00\)) already lie below, and within
1983--1995 the below-set is \(\{1983, 1984, 1986, 1990, \dots, 1995\}\)
--- nine of thirteen readings: the reference point does not separate
the eras. For this particular window three of seven readings fall below its
mean (the maximum is \(6.4\%\) above and the minimum \(5.5\%\)
below, outward-rounded). In general, the mean of nonconstant readings
falls between the minimum and maximum but guarantees neither a near-half
count below nor a uniform headroom bound.
\end{proposition}

\begin{proof}
All comparisons are direct rational comparisons of the printed
readings against \(276\) and \(\tfrac{8846}{10}\); the below-sets
follow by exact membership. The mean-rule statement is
Proposition~\ref{prop:windows}'s exact mean identity: a threshold at
the mean of a set of readings lies between the set's minimum and
maximum, so the set straddles it.
\end{proof}

\begin{proposition}[conditional stress bound: a sustained net SSB multiplier crosses the floor in sixteen steps]\label{prop:sustained}
Under the reference window's worst observed net SSB multiplier
\(\rho_{\mathrm{obs}} = \tfrac{41800}{45147}\), hypothetically held fixed
from the 1989 assessment reading \(921.66\) kt, the assessment-series
trajectory would fall below \(B_{\mathrm{aux}} = 276\) kt within sixteen
steps: \(\rho_{\mathrm{obs}}^{16} \le \tfrac{27600}{92166} <
\rho_{\mathrm{obs}}^{15}\). The count is a ceiling convention
(\(15.66\) steps rounds up). This is a stress calculation on a realized
net SSB ratio, not a prediction or a statement about the non-fishing
factor; the \(\kappa^*\) thresholds in Proposition~\ref{prop:brackets}
do not identify a sustained biological rate.
\end{proposition}

\begin{proof}
The sixteenth and fifteenth integer powers of
\(\tfrac{41800}{45147}\) are compared against
\(\tfrac{27600}{92166}\) by exact cross-multiplication. The ceiling
convention is stated because the exact threshold step lies strictly
between the fifteenth and sixteenth powers.
\end{proof}

\section{The outcome years}\label{outcome}

\begin{proposition}[assessment-series re-crossings and the conditional recovery-side threshold]
\label{prop:recovery}
Each SSB series re-crosses \(B_{\mathrm{aux}}=276\) kt at its 2015 point reading: the input vintage goes \(250.12\to298.65\) kt (2014--2015), and the database vintage \(238\to277\) kt. The database vintage then reads \(340,433,394,419,440,411\) kt for 2016--2021, all above the floor; its maximum, \(440\) kt, is \(\tfrac{44000}{88460}\approx49.7\%\) of \(B_{\mathrm{ref}}\). DFO 2025 Table~1 reports total landings of 4.4--12.9 kt per year for 2015--2021 (maximum 12.881 kt in 2017). The ratios of those reported landings to database-vintage SSB are 2.3--3.0\% for 2016--2021 and 1.60\% for 2015; the corresponding SSB-equivalent shares would be \(\kappa\) times these ratios, and are not identified here.

On the database-vintage 2020--2021 step, SSB falls from \(440\) to \(411\) kt, a point-series decline of \(29\) kt. Both adjacent DFO landings rows are shown: \(C_{2020}^{\rm rep}=10.260\) kt and \(C_{2021}^{\rm rep}=10.977\) kt. The resulting thresholds are \(\kappa^*_{2020,2020}=29/10.260\approx2.8265\) and \(\kappa^*_{2020,2021}=29/10.977\approx2.6419\). Under either accounting placement and both annual-row assignments, a subunitary accounting factor follows only if \(0\le\kappa<2.6419\); since \(\kappa\) is unknown, no recovery-side non-fishing attribution is established. The asymmetry described here is between the recorded assessment/database SSB trajectories, not a causal explanation of collapse or recovery.
\end{proposition}

\begin{proof}
The threshold crossings and six outcome values are direct rational comparisons in the two locked biomass files; the maximum comparison is the reduced ratio \(\tfrac{44000}{88460}\). The 2015--2021 landings are DFO Table~1 rows converted from tonnes to kt. The 2020--2021 thresholds divide the 29-kt database-vintage SSB decline by each adjacent-year reported-landings value; they are conditional on the same \(\kappa\) convention as Table~\ref{tab:brackets}.
\end{proof}

\begin{remark}[the monitoring consequence]\label{rem:monitoring}
Read together: the assessment input-vintage SSB point series contains
the collapse-era declines, while Propositions~\ref{prop:brackets} and
\ref{prop:moratorium} provide only conditional \(\kappa^*\) thresholds
for relating those declines to a reported-landings series. The raw
survey record also declines (Proposition~\ref{prop:survey}); the
capelin and mortality diagnostics cross in the same annual step
(Proposition~\ref{prop:capelin}); and the mean-rule reference point
separates nothing (Proposition~\ref{prop:discrimination}). None of
these series comparisons identifies a causal mechanism. The assessment
machinery itself was reconstructed around a monitoring indicator after
2016 (a three-year review cycle with an interim survey indicator, the
indicator value available by early January) --- the same review-timing
reading the companion paper certifies for the 2022--2023 observation
suspension (Abaee, 2026, The Northern cod forecast ladder, Section 3.10).

\begin{figure}[t]
\centering
\includegraphics[width=0.96\linewidth]{figs_arv/fig_record.pdf}
\caption{The certified record, both vintages: the input vintage
(NCAM 2016, 1983--2015) and the database vintage (RAM v4.66, 2014--2021)
against the auxiliary floor ($276$ kt) and the reference point
($884.6$ kt); reference and collapse windows shaded; the 1993 breach
($101.05$ kt), the 1995 minimum ($9.68$ kt), the 22 consecutive
below-floor readings to 2014, and the 2015 re-crossing
($250.12 \to 298.65$; $238 \to 277$) as certified in the text.
Generated by the assertion script
\texttt{paper2\_arv\_record\_figure\_v1.py} (recompute-then-assert;
eleven exact checks against the locked files).}
\label{fig:record}
\end{figure}
\end{remark}

\begin{table*}[t]
\centering
\scriptsize
\setlength{\tabcolsep}{2.8pt}
\caption{Selected recorded values and conditional thresholds. SSB is in kt; DFO 2025 Table~1 landings are reported in tonnes unless noted. Thresholds are dimensionless \(\kappa^*\) values using reported landings converted to kt; displayed values are rounded.}
\label{tab:summary}
\begin{tabular}{@{}>{\raggedright\arraybackslash}p{.36\textwidth}>{\raggedright\arraybackslash}p{.37\textwidth}>{\raggedright\arraybackslash}p{.22\textwidth}@{}}
\toprule
Quantity & Value & Source file \\
\midrule
Ref-window worst net ratio (1985--1986) & \(41800/45147\) & Table A2 \\
Collapse net ratios (1991--1992; 1992--1993) & \(0.520,\ 0.265\) & Table A2 \\
Collapse \(\kappa^*\) pairs \((C_t,C_{t+1})\) & \(\begin{gathered}2.0496,\ 8.6083;\\6.8586,\ 24.6577\end{gathered}\) & A2 + T1 \\
Post-collapse \(\kappa^*\) pairs & \(\begin{gathered}6.1886,\ 53.6530;\\15.8828,\ 50.5327\end{gathered}\) & A2 + T1 \\
Post-collapse two-year SSB ratio & \(10105/968 \approx 10.44\) & Table A2 \\
Reported landings, 1992--1993 & \(40{,}956,\ 11{,}392\) t & Table 1 (2025) \\
Reported landings, 1994--1995 & \(1{,}314,\ 413\) t & T1; catch reconstruction \\
2020--2021 \(C_t,C_{t+1}\) & \(10.260,\ 10.977\) kt & Table 1 (2025) \\
2020--2021 \(\kappa^*_{t,t},\kappa^*_{t,t+1}\) & \(2.8265,\ 2.6419\) & DB SSB + T1 \\
Survey index, 1989/1994 & \(2{,}127{,}417,\ 21{,}797\) & Survey T3 \\
Capelin, 1990/1991 & \(5783,\ 138\) kt & Capelin file \\
First below-floor assessment point & 1993 (\(101.05 < 276\)) & Table A2 \\
Cross-vintage 1993 SSB & \(101.05\) vs \(79\) kt & A2 + T17 \\
Outcome years 2016--2021 & \(340\)--\(440\) kt & DB extract \\
\bottomrule
\end{tabular}
\end{table*}

\section{Delimitations}\label{scope}

Every arithmetic statement concerns the printed rows of the locked
files, not a forecast or a stochastic model. The assessed biomass
series are model-based SSB outputs, and DFO 2025 Table~1 reports total
landings rather than all fishing removals. The unit conversion from
tonnes to kilotonnes is exact, but \(\kappa\), the conversion from
reported landings/catch to SSB-equivalent removals, is unknown and is
not estimated. Lemma~\ref{lem:bracket} therefore gives only a
conditional accounting implication: for a selected row and either
placement, \(g_t<1\) follows only if \(\kappa<\kappa^*\); robust use
across adjacent-year assignments requires satisfying both thresholds.
The result is neither an identified zero-removal counterfactual nor a
biological or causal conclusion. The published uncertainty bands are
left untouched; exact arithmetic on point estimates does not remove
assessment uncertainty or establish the latent stock trajectory. The
mortality diagnostic is model-estimated given assumed catch, so its
crossing does not independently identify a mechanism. The survey index
has its own sampling error and catchability assumptions; the
instrument--output divergence is a difference between recorded series,
not a causal attribution. The capelin series has observed years only,
and the contemporaneity statement is at annual resolution. The scope
remains one stock, two vintages for the auxiliary threshold, and one
input-vintage reference point constructed as its own seven-year mean.

\section{Verification methods}\label{methods}

The companion arithmetic checker,
\texttt{verify\_kappa\_thresholds.py}, uses only the Python
standard library and reads the five small CSV extracts in
\texttt{source\_data/}. It checks the selected NCAM 2016 Table~A2 SSB
rows, the RAM Legacy database-vintage SSB rows for 2020--2021, and the
DFO 2025 Table~1 reported-total-landings rows; it also checks the
1992--1995 and 2015 cross-file equalities against the Schijns catch
reconstruction. The checker verifies the seven-row 1983--1989
landings-range calculation, scans all seventy consecutive xteNCAM
assessment-band pairs, checks the displayed band endpoints and floor
comparisons, converts tonnes to kilotonnes exactly, recomputes all seven
primary \(\kappa^*\) pairs and the separate 2020--2021 pair in rational
arithmetic, checks the displayed values, and tests both accounting
forms below, at, and above each threshold. This is a source-extract
consistency and arithmetic check, not a re-extraction of the cited
reports: it does not validate the underlying assessment estimates or
their uncertainty, estimate \(\kappa\), resolve within-year timing, or
cover the other survey, mortality, and covariate summaries.

\section{Conclusion}\label{conclusion}

The principal result is a conditional accounting boundary, not an
estimate of removals or a biological attribution. NCAM assessment SSB
falls by \(352.56\) kt from 1991 to 1992 and by \(280.90\) kt from
1992 to 1993. Against DFO 2025 Table~1's reported total landings, the
row-specific thresholds for these two transitions are respectively
\((2.0496,8.6083)\) and \((6.8586,24.6577)\), where each pair uses
\(C_t^{\rm rep}\) and \(C_{t+1}^{\rm rep}\). Consequently, for both
transitions, both adjacent-year assignments, and either accounting
placement, the conditional multiplier is below one only when
\(0\leq\kappa<2.0496\). The selected 1993--1995 steps have a minimum
threshold of \(6.1886\) across the same placements. For the separate
2020--2021 RAM database-vintage point series, the \(29\) kt decline and
reported landings of \(10.260\) and \(10.977\) kt yield thresholds
\(2.8265\) and \(2.6419\); a subunitary multiplier across both row
assignments follows only for \(0\leq\kappa<2.6419\).

These thresholds preserve rather than conceal two unresolved inputs:
\(\kappa\) is unknown, and the sources do not fix a unique within-year
alignment between SSB transitions and reported landings. DFO's total-
landings series is not relabelled as all removals. The method's
contribution is therefore a reproducible, two-placement boundary on
what the displayed series can imply, not an unconditional contraction
certificate, an estimate of a harvest-free trajectory, or a causal
claim. The survey-index decline, the annual coincidence of the
capelin and mortality diagnostics, and the assessment-vintage changes
remain descriptive cross-series evidence; they do not identify a
mechanism. All thresholds use published point estimates and do not
supplant the sources' uncertainty or resolve assessment-model
assumptions. The \(276\) kt floor and \(884.6\) kt reference point are
likewise auxiliary conventions for these comparisons, not universal
biological boundaries.

\section*{References}

{\footnotesize
DFO (2016). Stock assessment of Northern cod (NAFO 2J3KL). \emph{Can.
Sci. Advis. Sec. Sci. Advis. Rep.} 2016/026.

Regular, P. et al. (2025). Northern cod (NAFO 2J3KL) reassessment.
\emph{DFO Can. Sci. Advis. Sec. Res. Doc.} 2025/048.

Ricard, D. et al. (2012). Examining the knowledge base and status of
commercially exploited marine species with the RAM Legacy Stock
Assessment Database. \emph{Fish and Fisheries} 13, 380--398.

Schijns, R. et al. (2021). Reconstructing three centuries of fish and
fisheries catch in the Northwest Atlantic. \emph{ICES Journal of Marine
Science} 78, 2675--2684.

Abaee, A. (2026). An Obstruction Calculus for Viability under Incomplete Observation. Figshare preprint, version 2. \url{https://doi.org/10.6084/m9.figshare.33716593.v2}.

Abaee, A. (2026). The Northern cod forecast ladder under a frozen
specification. Preprint.

Abaee, A. (2026). Calibration data record, Northern cod (NAFO 2J3KL):
sources, cross-checks, and locked extracts. Preprint.

Hutchings, J.A., Myers, R.A. (1994). What can be learned from the
collapse of a renewable resource? Atlantic cod, Gadus morhua, of
Newfoundland and Labrador. \emph{Canadian Journal of Fisheries and
Aquatic Sciences} 51, 2126--2146.

Aubin, J.-P. (1991). Viability Theory. Birkh\"auser, Boston.

Myers, R.A. and Cadigan, N.G. (1995). Was an increase in natural mortality responsible for
the collapse of northern cod? \emph{Can. J. Fish. Aquat. Sci.} \textbf{52}, 1274--1285.

Myers, R.A., Hutchings, J.A. and Barrowman, N.J. (1997). Why do fish stocks collapse? The
example of cod in Atlantic Canada. \emph{Ecol. Appl.} \textbf{7}, 91--106.
}

\section*{Declarations}

\section*{Funding}
None.

\section*{Competing interests}
None.

\section*{Data availability}
All inputs are the archived public files of the calibration record
(wave-e-cod extracts and the record's CSV files), with primary sources
cited at archival time.

\section*{Code availability}
The version-specific checker
\texttt{verify\_kappa\_thresholds.py} and its \texttt{source\_data/}
inputs accompany this edition in the versioned
\texttt{paper09b\_arv\_certification\_v6\_src/} folder. It verifies
the source-row checks and conditional-threshold arithmetic described
above; it is not a full re-execution of every descriptive result in the
paper. The inherited
record figure is generated by
\texttt{paper2\_arv\_record\_figure\_v1.py}, available in the repository's
\texttt{arena agent 1/paper rewrites/latex/} folder; its original assertions
are separate from this threshold checker.

\section*{AI declaration}
GLM (Z.ai), Qwen (Alibaba Cloud) and DeepSeek AI assisted with drafting
and iterative review. The author reviewed and edited outputs and takes
responsibility for the final work.

\end{document}